\documentclass{article}

\usepackage[preprint]{neurips_2023}

\usepackage{amsmath,amssymb,amsfonts}
\usepackage{amsthm}
\usepackage{graphicx}
\usepackage{booktabs}
\usepackage{hyperref}
\usepackage{url}
\usepackage{algorithm}
\usepackage{algorithmic}
\usepackage{xcolor}
\usepackage{subcaption}
\usepackage{microtype}
\usepackage[capitalize,noabbrev]{cleveref}

\PassOptionsToPackage{numbers,sort&compress}{natbib}
\graphicspath{{figs/}}

\newcommand{\E}{\mathbb{E}}
\newcommand{\R}{\mathbb{R}}
\newcommand{\norm}[1]{\left\lVert #1 \right\rVert}

\title{Dendrites vs Backprop: Local Three-Factor Learning on Nonlinear Tasks}

\author{
  name name \\
    University \\
  \texttt{blah}
}

\begin{document}
\maketitle

\begin{abstract}
% TODO: Write a short summary of the project motivation, experiments, and key findings.
\end{abstract}

% ---------------------------
\section{Introduction}
% - Motivation: biological plausibility vs backprop
% - Core question: how far can a single neuron with nonlinear dendrites go?
% - Contributions summary

% ---------------------------
\section{Background and Related Work}
% - Dendritic computation and compartmental models
% - Local three-factor rules (pre, post, modulator)
% - Biological vs artificial credit assignment
% - Prior comparisons with backprop

% ---------------------------
\section{Model and Learning Rules}
% \subsection{Two-compartment neuron model}
% \subsection{Three-factor local learning rule}
% \subsection{Eligibility traces and stability}
% \subsection{Equivalent backprop baseline}
% \subsection{Tiny MLP baseline}

% ---------------------------
\section{Tasks and Evaluation}
% \subsection{Datasets and tasks (XOR, 3-bit parity, Boolean families)}
% \subsection{Metrics: accuracy, sample efficiency, robustness}
% \subsection{Noise robustness tests}

% ---------------------------
\section{Training Algorithms}
% \subsection{Local three-factor algorithm (pseudo-code)}
% \subsection{Backprop training procedure}
% \subsection{Implementation details (optimizer, learning rate, epochs)}

% ---------------------------
\section{Experiments}
% \subsection{Setup and hyperparameters}
% \subsection{Results: comparison tables and learning curves}
% \subsection{Ablations: eligibility traces, nonlinearity choice}
% \subsection{Analysis: sample efficiency and noise robustness}

% ---------------------------
\section{Discussion}
% - Interpretation of results
% - When dendritic models outperform / fail
% - Trade-offs: stability vs sample efficiency

% ---------------------------
\section{Conclusion}
% - Summary of findings
% - Implications for biologically plausible learning
% - Future work directions

% ---------------------------
\section{Limitations}
% - Scope, assumptions, and dataset scale

% ---------------------------
\section{Ethical Considerations}
% - No sensitive data used, educational research context

% ---------------------------
\bibliographystyle{plainnat}
\bibliography{refs}

\appendix
\section{Appendix}
% - Additional figures, code listings, and hyperparameters

\end{document}
